\section{用赝PDF方法由格点计算得到胶子PDF}\label{sec:cal-details}

本工作中，我们使用文献~\cite{Balitsky:2019krf}定义的非极化胶子算符，
%
\begin{equation}\label{eq:gluon_operator}
 {\cal O(z)}\equiv\sum_{i\neq z,t}{\cal O}(F^{ti},F^{ti};z)-\sum_{i,j\neq z,t}{\cal O}(F^{ij},F^{ij};z),
\end{equation}
%
其中算符${\cal O}(F^{\mu\nu}, F^{\alpha\beta};z) = F^\mu_\nu(z)U(z,0)F^{\alpha}_{\beta}(0)$，$z$为 Wilson 线长度，场强$F_{\mu\nu}$定义为
%
\begin{equation}\label{}
 F_{\mu\nu} = \frac{i}{8a^2g_0}\left(\mathcal{P}_{[\mu,\nu]}+\mathcal{P}_{[\nu,-\mu]}+\mathcal{P}_{[-\mu,-\nu]}+\mathcal{P}_{[-\nu,\mu]}\right),
\end{equation}
%
其中$a$是晶格间距，
$g_0$是强耦合常数，
plaquette $\mathcal{P_{\mu,\nu}}=U_{\mu}(x)U_{\nu}(x+a\hat{\mu})U^{\dag}_{\mu}(x+a\hat{\nu})U^{\dag}_{\nu}(x)$，而$\mathcal{P_{[\mu,\nu]}}=\mathcal{P}_{\mu,\nu}-\mathcal{P}_{\nu,\mu}$。
还存在一个可选算符$\sum_{i\neq z,t}{\cal O}(F^{ti},F^{zi};z)$，对应文献~\cite{Balitsky:2019krf}中相同的匹配核。
我们不选该算符，因为它因运动学原因在$P_z=0$时为零，会给获取分布带来额外困难。

利用式~\ref{eq:gluon_operator}的算符，我们计算基态介子$|0 (P_z) \rangle$在不同提升动量$P_z$与 Wilson 线位移长度$z$下的格点胶子矩阵元，
随后研究它们对Ioffe时间$\nu=zP_z$的依赖关系
%
\begin{equation}
\mathcal{M}(\nu,z^2) = \langle 0 (P_z)|{\cal O}(z)|0 (P_z)\rangle,
\label{eq:ME_unpol}
\end{equation}
%
并称$\mathcal{M}(\nu,z^2)$为Ioffe时间分布（ITD）。
为消除ITD中的紫外发散，我们把ITD除以$P_z=0$处相应的$z$依赖矩阵元来构造约化ITD（RITD），再如首个夸克赝PDF计算~\cite{Orginos:2017kos}那样用$z^2=0$处的矩阵元对该比值作归一化，
%
\begin{equation}
\mathscr{M}(\nu,z^2)=
\frac{\mathcal{M}(zP_z,z^2)/\mathcal{M}(0\cdot P_z,0)}{\mathcal{M}(z\cdot 0,z^2)/\mathcal{M}(0\cdot 0,0)}.
\label{eq:RITD}
\end{equation}
%
${\cal O}(z)$的重正化与运动学因子在RITD中相互抵消。
这里使用的RITD双比值在$z=0$处自动归一到1。

RITD通过赝PDF匹配条件~\cite{Balitsky:2019krf}与$\pi$介子的胶子分布$g$及夸克分布$q_S$相联系：
%
\begin{align}
\mathscr{M}(\nu,z^2)&=\int_0^1 dx \frac{xg(x,\mu^2)}{\langle x \rangle_g}R_{gg}(x\nu,z^2\mu^2) \nonumber \\
&+\frac{P_z}{P_0}\int_0^1 dx \frac{xq_S(x,\mu^2)}{\langle x \rangle_g}R_{gq}(x\nu,z^2\mu^2),
\label{matching-eq}
\end{align}
%
其中$\mu$是$\overline{\text{MS}}$方案的重正化能标，
$\langle x \rangle_g=\int_0^1 dx \, x g(x,\mu^2)$为胶子在$\pi$介子中的动量份额。
我们可以把式~\ref{matching-eq}中胶子内嵌于胶子的$R_{gg}$贡献拆成两部分，即文献~\cite{Radyushkin:2018cvn}中式21与式24引入的两项：
%
\begin{align}
&R_{gg}(y,z^2\mu^2)=R_1(y,z^2\mu^2)+R_2(y), \label{matching-kernel}\\
&R_1(y,z^2\mu^2)=-\frac{\alpha_s(\mu)}{2\pi}N_c\ln\left(z^2\mu^2\frac{e^{2\gamma_E+1}}{4}\right)R_B(y), \label{matching-evolve}\\
&R_2(y)=\cos y-\frac{\alpha_s(\mu)}{2\pi}N_c\left( 2R_B(y)+R_L(y)+R_C(y)\right),
\label{matching-scheme}
\end{align}
%
其中$R_1(y,z^2\mu^2)$是与演化有关的项，
$R_2(y)$
是与方案转换有关的项，
$\alpha_s$为能标$\mu$处的强耦合常数，
$N_c=3$是色数，
$\gamma_E=0.5772$是欧拉--马歇罗尼常数。
%
把$R_1(y,z^2\mu^2)$中的$z$取为$2e^{-\gamma_E-1/2}/\mu$可使对数项为零，从而压低含对数项高阶的残余，正如赝PDF一圈演化的论文~\cite{Radyushkin:2018cvn}所讨论的那样。
%
胶子内嵌于夸克的核$R_{gq}(y,z^2\mu^2)$以及$R_B(y)$、$R_L(y)$与$R_C(y)$定义于文献~\cite{Balitsky:2019krf}的式7.21--23及其下一段。

本工作中我们首先忽略$\pi$介子夸克PDF，因为在全局拟合中总夸克PDF被发现远小于胶子PDF~\cite{Barry:2018ort,Novikov:2020snp}；
稍后我们会估计这一假设引入的系统不确定度。
演化后的胶子ITD（EITD）$G$由演化项$R_1(y,z^2\mu^2)$得到：
%
\begin{align}
G(\nu,\mu,z^2)&=\mathscr{M}(\nu,z^2)\nonumber\\
&+\int^1_0 dx\, R_1(x,z^2\mu^2)\mathscr{M}(x\nu,z^2).
\label{evolution}
\end{align}
%
EITD的$z$依赖应被演化公式中的$\ln{z^2}$项所补偿。
原则上EITD $G$不含$z$依赖，并通过方案转换项$R_2(y)$与光锥胶子PDF $g(x,\mu^2)$相连：
%
\begin{equation}
G(\nu,\mu)=\int_0^1 dx\, \frac{xg(x,\mu^2)}{\langle x \rangle_g} R_2(x\nu),
\label{conversion}
\end{equation}
%
因此反解该方程即可提取胶子PDF $g(x,\mu^2)$。

在格点上，我们在三个系综上使用 clover 价费米子；这些系综具有 MILC 合作组~\cite{Bazavov:2012xda}产生的$N_f = 2+1+1$味高度改进 staggered 夸克（HISQ）~\cite{Follana:2006rc}，涵盖两种晶格间距（$a\approx 0.12$与0.15~fm）和三种$\pi$介子质量（220、310、690~MeV）。
clover 夸克的质量被调到重现 PNDME 合作组~\cite{Rajan:2017lxk,Bhattacharya:2015wna,Bhattacharya:2015esa,Bhattacharya:2013ehc}所用最轻的海轻、海奇异赝标介子质量。
%
我们在胶子圈上施加五步 HYP 涂抹~\cite{Hasenfratz:2001hp}以减小统计不确定度，如文献~\cite{Fan:2018dxu}中所研究。
我们对夸克场使用高斯动量涂抹~\cite{Bali:2016lva}以达到更高的介子提升动量。
表~\ref{table-data}列出三个系综的晶格间距$a$、价$\pi$介子质量$M_\pi^\text{val}$与$\eta_s$质量$M_{\eta_s}^\text{val}$、格点大小$L^3\times T$、组型数$N_\text{cfg}$、两点关联函数总测量数$N_\text{meas}^\text{2pt}$，以及三点关联函数拟合所用的时间分离$t_\text{sep}$。
这使我们能够通过外推到达连续极限与物理$\pi$介子质量。
不同系综的总测量数在$10^5$--$10^6$之间变化。


\begin{table}[!htbp]
\centering
\begin{tabular}{|c|c|c|c|}
\hline
  系综  & a12m220 & a12m310 & a15m310 \\
\hline
  $a$ (fm)  & $0.1184(10)$ & $0.1207(11)$  & $0.1510(20)$ \\
\hline
  $M_{\pi}^\text{val}$ (MeV)  & $226.6(3)$ & $311.1(6)$ & $319.1(31)$\\
\hline
  $M_{\eta_s}^\text{val}$ (MeV)  & $696.9(2)$ & $684.1(6)$ & $687.3(13)$\\
\hline
  $L^3\times T$  & $32^3\times 64$ & $24^3\times 64$ & $16^3\times 48$ \\
\hline
  $P_z$ (GeV)  & $[0,2.29]$ &  $[0,2.14]$ &  $[0,2.05]$ \\
\hline
  $N_\text{cfg}$  & 957 & 1013 & 900 \\
\hline
  $N_\text{meas}^\text{2pt}$  &  731,200  & 324,160 & 21,600 \\
\hline
  $t_\text{sep}$  & \{5,6,7,8,9\}   & \{5,6,7,8,9\} & \{4,5,6,7\} \\
\hline
\end{tabular}
\caption{
本研究分析的三个系综——由 MILC 合作组产生的$N_f=2+1+1$ HISQ 背景上的 clover 价费米子——所用的晶格间距$a$、价$\pi$介子质量$M_\pi^\text{val}$与$\eta_s$质量$M_{\eta_s}^\text{val}$、格点大小$L^3\times T$、组型数$N_\text{cfg}$、两点关联函数总测量数$N_\text{meas}^\text{2pt}$以及三点关联函数拟合的时间分离$t_\text{sep}$。
}
\label{table-data}
\end{table}


介子$\Phi$的两点关联函数为
%
\begin{align}\label{}
C_\Phi^\text{2pt}(P_z;t)&=
 \int dy^3 e^{-i y\cdot P_z} \langle \chi_\Phi(\vec{y},t)|\chi_\Phi(\vec{0},0)\rangle \nonumber \\
  &= |A_{\Phi,0}|^2 e^{-E_{\Phi,0}t} + |A_{\Phi,1}|^2 e^{-E_{\Phi,1}t} + ...,
\label{eq:2pt_fit_formula}
\end{align}
%
其中$P_z$是介子沿$z$方向的动量，$\chi_\Phi=\bar{q}_1\gamma_5 q_2$为赝标介子插值算符，
$t$是欧氏时间，
而$|A_{\Phi,i}|^2$与$E_{\Phi,i}$分别是基态（$i=0$）和第一激发态（$i=1$）的振幅与能量。


\begin{figure}[htbp]
\centering
	\centering
\includegraphics[width=0.9\textwidth]{images/Ratio_c957_O0IB_dir-average_p2_z1_Fits.pdf}
\includegraphics[width=0.9\textwidth]{images/Ratio_c957_O0IB_dir-average_p4_z4_Fits.pdf}
\caption{
来自最轻$\pi$介子质量系综（$a\approx 0.12$~fm、$M_\pi\approx 220$~MeV）在$P_z=2\times 2\pi/L$、$z=1$（上行）与$P_z=4\times 2\pi/L$、$z=4$（下行）的比值图示例（左列）、单态拟合（第二列）与 two-sim 拟合（后两列）。
所有图中灰色带是用$t_\text{sep}\in[5,9]$的 two-sim 拟合提取的基态矩阵元。
从左到右各列为：
三点与两点关联函数之比及其由$t_\text{sep}\in [5,9]$ two-sim 拟合给出的重构拟合带（表示为$t-t_\text{sep}/2$的函数）；
各$t_\text{sep}\in[3,9]$处三点关联函数的单态拟合结果；
使用$t_\text{sep}\in[t_\text{sep}^\text{min},9]$的 two-sim 拟合结果随$t_\text{sep}^\text{min}$的变化；以及
使用$t_\text{sep}\in[5,t_\text{sep}^\text{max}]$
的 two-sim 拟合结果随$t_\text{sep}^\text{max}$的变化。}
\label{fig:Ratio-fitcomp}
\end{figure}


三点胶子关联函数由胶子圈与$\pi$介子两点关联函数组合得到。
把三点关联函数拟合到能量本征态展开即可获得胶子算符的矩阵元：
%
\begin{align}\label{eq:3ptC}
&C_\Phi^\text{3pt}(z,P_z;t_\text{sep},t) \nonumber \\
&=\int d^3y\, e^{-iy\cdot P_z}\langle \chi_\Phi(\vec y,t_\text{sep})|{\cal O}(z,t)|\chi_\Phi(\vec 0,0)\rangle \nonumber \\
&= |A_{\Phi,0}|^2\langle 0|{\cal O}|0\rangle e^{-E_{\Phi,0}t_\text{sep}} \nonumber \\
&+ |A_{\Phi,0}||A_{\Phi,1}|\langle 0|{\cal O}|1\rangle e^{-E_{\Phi,1}(t_\text{sep}-t)}e^{-E_{\Phi,0}t} \nonumber \\
&+ |A_{\Phi,0}||A_{\Phi,1}|\langle 1|{\cal O}|0\rangle e^{-E_{\Phi,0}(t_\text{sep}-t)}e^{-E_{\Phi,1}t} \nonumber \\
&+ |A_{\Phi,1}|^2\langle 1|{\cal O}|1\rangle e^{-E_{\Phi,1}t_\text{sep}}+ ...,
\end{align}
%
其中$t_\text{sep}$为源--汇时间分离，
$t$是胶子算符插入时间。
振幅与能量$A_{\Phi,0}$、$A_{\Phi,1}$、$E_{\Phi,0}$与$E_{\Phi,1}$由两点关联函数的两态拟合得到。
$\langle 0|{\cal O}|0\rangle$、$\langle 0|{\cal O}|1\rangle$（$\langle 1|{\cal O}|0\rangle$）与$\langle 1|{\cal O}|1\rangle$分别是基态矩阵元、基态--激发态矩阵元与激发态矩阵元。
我们从三点关联函数的两态拟合提取基态矩阵元$\langle 0|{\cal O}|0\rangle$，或在多个时间分离上对含$\langle 0|{\cal O}|0\rangle$、$\langle 0|{\cal O}|1\rangle$与$\langle 1|{\cal O}|0\rangle$项作两态同时（``two-sim''）拟合。

为验证拟合得到的矩阵元可靠，我们将其与三点、两点关联函数之比
%
\begin{equation}\label{eq:Ratio}
R^\text{ratio}(z,P_z;t_\text{sep},t) = \frac{C^\text{3pt}(z,P_z;t_\text{sep},t)}{C^\text{2pt}(P_z;t_\text{sep})};
\end{equation}
%
比较：若无激发态贡献，该比值应等于基态矩阵元。
图~\ref{fig:Ratio-fitcomp}左列展示了最轻$\pi$介子系综 a12m220 在选定动量$P_z$与 Wilson 线长度$z$下胶子矩阵元比值的示例。
我们看到比值随源--汇分离从0.60增大到1.08~fm而上升；
在大分离处比值开始收敛，表明忽略激发态带来的问题减轻。
灰带是对五个$t_\text{sep}$的三点关联函数作 two-sim 拟合得到的基态矩阵元。
忽略激发态的拟合之收敛性也可见图~\ref{fig:Ratio-fitcomp}第二列，其中比较了各源--汇分离的单态拟合结果：
单态拟合结果随$t_\text{sep}$增大而上升，在大$t_\text{sep}$处开始收敛到 two-sim 拟合结果。

图~\ref{fig:Ratio-fitcomp}第三、四列展示分别使用$t_\text{sep}\in[t_\text{sep}^\text{min},9]$
与
$t_\text{sep}\in[5,t_\text{sep}^\text{max}]$
的 two-sim 拟合，以研究 two-sim 基态矩阵元对纳入拟合的源--汇分离范围的依赖。
我们观察到各矩阵元在一个标准差内彼此一致，表明基态矩阵元提取一致，尽管统计误差大于单态拟合。
当采用较小的$t_\text{sep}^\text{min}=3$与4、或较小的$t_\text{sep}^\text{max}=6$与7时，矩阵元提取的涨落较大；当$t_\text{sep}^\text{min}>4$且$t_\text{sep}^\text{max}>7$时，two-sim 拟合给出的基态矩阵元非常稳定。

图~\ref{fig:RITD-fitcomp}展示了相同例子（$P_z=2\times 2\pi/L$、$z=1$与$P_z=4\times 2\pi/L$、$z=4$）用$t_\text{sep}\in[t_\text{sep}^\text{min},9]$的 two-sim 拟合结果得到的RITD。
为压制格点涨落而构造的RITD结果在所考虑的各种拟合下都非常稳定。
对于 a12m310 与 a15m310 系综，由于$\pi$介子更重，其RITD的$t_\text{sep}$依赖比 a12m220 系综更平缓。
总体而言，我们从 two-sim 拟合得到的基态RITD是稳定的，并用它们提取胶子PDF。


\begin{figure}[htbp]
\centering
	\centering
	\includegraphics[width=0.43\textwidth]{images/RITD_c957_O0IB_p2_z1_tsepmin.pdf}
	\includegraphics[width=0.43\textwidth]{images/RITD_c957_O0IB_p4_z4_tsepmin.pdf}
\caption{
来自 a12m220 系综的RITD示例，表示为$t_\text{sep}^\text{min}$的函数：上图为$P_z=2\times 2\pi/L$、$z=1$，下图为$P_z=4\times 2\pi/L$、$z=4$。
使用$t_\text{sep}\in[t_\text{sep}^\text{min},9]$的 two-sim 拟合RITD结果与最终选用$t_\text{sep}\in[5,9]$的结果一致。}
\label{fig:RITD-fitcomp}
\end{figure}

